% Transcription — fonds Alexandre Grothendieck, shelfmark 69, batch 3 (pages 41-60).
% Established from the facsimile archives/batches/69/batch-03.pdf (cut from a
% copy of 69.pdf fetched through the project's relay,
% https://grothendieck-relay.onrender.com/source/69.pdf, a mirror of
% https://grothendieck.umontpellier.fr/69.pdf), per the skill
% transcribe-grothendieck. The raw scan's cover sheet is not in the batch
% file: PDF page k is archivists' page 40+k, and the pencilled numbers
% bottom-left (« 41 » ... « 60 ») agree. The folder is 119 pages in six
% batches (1-20, 21-40, 41-60, 61-80, 81-100, 101-119), transcribed in
% parallel; this file does not rely on its neighbours.
%
% Pass: Opus 5.5 (claude-opus-5-5), 2026-09-27 — first pass, unchecked against the pages by a human.
%
% Pages skipped: none. Pages summarised: none. All twenty pages are his
% working notes in blue ink on punched listing paper (staves ruled in grey),
% and all are transcribed.
%
% His pagination. Page 41 opens a run headed « Les 27 droites de la surface
% cubique, et le syst. de racines E6 », paginated by him 1 to 17 at the top
% of pages 41-57 (« 1 » over the title, 2-17 circled except 7 and 8). Pages
% 58-60 carry no number of his and restart under the underlined heading
% « Graphes cubiques ». Internal references (« (12) à (16) de (I) ») refer to
% his numbered formulas within the run.
%
% What the batch is about. The combinatorics of the 27 lines on a cubic
% surface and the root system E6, which is where the folder's « graphes
% cubiques » come from. I (pp. 41-43): P blown up at six general points
% s_1..s_6; the curves D_i, D'_i (conics), D_ij (lines), their classes
% eta, xi_i, xi'_i, xi_ij in Pic, the relations (1)-(16), lambda = 3 eta - xi
% (the cubic embedding, K = -3 lambda in a margin note), each line meeting
% ten others and any two meeting lines having a unique third. II (pp.
% 43-48): the same algebra over a ring k, the dual family (eta', xi'_i), the
% intersection form, the « graphe cubique » Delta on the 27 elements (linked
% iff product 1), the lemma on the unique third vertex, the definition of a
% cubic graph as a graph isomorphic to the standard one, the lattice
% E^(Delta) = k^[Delta]/N(Delta) of rank 7, ordered bases, Isom(Delta_0,
% Delta) a torsor, the canonical lambda_Delta with lambda.s = 1. III (pp.
% 49-57): roots as differences of unlinked vertices, the 72 roots in three
% families r_0, r_ij, r_ijk each tied to a « double six », the reflections
% s_r, W = Aut(Delta) transitive on roots and bases, Card W = 72.6! =
% 2^7 3^4 5, the theorem that R(Delta) is a root system of type E6 with
% simple roots r_1..r_6 and the Dynkin diagram (redrawn in tikz-cd), the
% highest root r_0, W of index 2 in Aut of the root system, the map
% i(delta) = 3 delta - lambda to vectors of square -12 and the splitting of
% Aut(R) as W x Z/2. Pp. 58-60 restart: an intrinsic definition of cubic
% graphs by conditions a)-e) on two sets I, I' of six vertices, bases and
% « bibases », the 72 bases listed in four families (i)-(iv), corollaries,
% and a cancelled attempt (crossed out on p. 60) at « bi-bases » and triples
% of linked vertices.
%
% Notation, fixed once. His gothic P is \mathfrak{P} (P_2(I), sets of pairs);
% his « E chapeau » is \widehat{E}; the symmetric group is \mathfrak{S}_6.
% His circled crossed reference marks are rendered (*) with a \note. Formulas
% where he wrote over earlier symbols are given as the final reading, and
% the overwriting is noted where the earlier symbol reads. Where a formula on
% the page disagrees with the computation around it (p. 44, p. 57), it is
% kept and flagged in a \note, not corrected. Pages 50 (lower half) and 60
% (lower half) are cancelled by long strokes and transcribed as far as they
% read, their many overwritings making them the hardest of the batch.
%
% DATING: the folder has its own date in the catalogue, copied verbatim; no
% group sentence.
\documentclass[11pt,a4paper]{article}
\input{../preamble/grothendieck}

\folder{69}
\batch{3}
\pages{41}{60}
\dating{[à partir de 1976]}
\watermark{Édition de démonstration}
\foldertitle{Graphes cubiques : notes manuscrites (s.d.), lettre (s.d.).}

\begin{document}

\page{41}
\section{Les 27 droites de la surface cubique, et le syst.\ de racines $E_6$}
\note{titre de sa main, en tête de la page ; au-dessus de « cubique », un
chiffre qui se lit « 1 », début de sa pagination de 1 à 17 (pages 41 à 57),
les suivants étant cerclés}

\textbf{I.} $P$ plan projectif (\uncertain{sur un corps} \ill{} \ill{} \ill{})

$S = \{s_1, \dots, s_6\}$ un ens.\ de $6$ pts de $P$ \struck{\ill{}}
\struck{trois \ill{} \ill{}} pas \add{trois à trois} alignés et pas
\struck{\ill{}} \add{tous} sur une conique

$\widetilde{P}$ déduit de $P$ en faisant éclater $S$

$\widetilde{S} = \sum_{1}^{6} D_i$ image inverse de $S$

$D_{ij}$ ($1 \leq i \neq j \leq 6$) image inv.\ propre de la droite qui joint
$s_i, s_j$

$D'_i$ image inverse propre de \emph{la} conique dans $P$ passant par
$S - \{s_i\}$
\marginal{vérifier}

On a
\[
(\ast)\quad
\begin{cases}
D_i \cap D_j = \emptyset & \text{si } i \neq j \\
D_i \cap D_{jk} \neq \emptyset & \text{ssi } i \in \{j, k\} \\
D'_i \cap D'_j = \emptyset & \text{si } i \neq j \\
D'_i \cap D_{jk} \neq \emptyset & \text{ssi } i \in \{j, k\} \\
D_i \cdot D'_j \neq \emptyset & \text{ssi } i \neq j \\
D_{ij} \cdot D_{kl} \neq \emptyset & \text{ssi } \{i,j\} \cap \{k,l\} \text{ de card.\ } 0 \text{ ou } 2
\end{cases}
\]
\note{le signe de renvoi devant l'accolade est un rond plein, rendu ici
$(\ast)$ ; la condition « ssi $\{i,j\} \cap \{k,l\}$ de card.\ $0$ ou $2$ » est
écrite au bout de l'avant-dernière ligne et vaut pour la dernière}
\marginal{vérifier — cf calculs plus bas}

On calcule, \emph{dans} \add{$E$} $\simeq A^1(\widetilde{P}) =
\operatorname{Pic}(\widetilde{P})$, en désignant par
\[
(1) \qquad \eta, \ (\xi_i)_{1 \leq i \leq 6}
\]
les images inverses de $\operatorname{cl}(\mathcal{O}_P(1))$ et les classes
des $D_i$, par $\xi'_i$ les classes des $D'_i$, par $\xi_{ij}$ celles des
$D_{ij}$. On trouve
\[
\begin{aligned}
&(2) & \xi_{ij} &= \eta - \xi_i - \xi_j \\
&(3) & \xi'_i &= 2\eta - \sum_{\alpha \neq i} \xi_\alpha = 2\eta - \xi + \xi_i
\end{aligned}
\]
i.e.
\[
(3\,\text{bis}) \qquad \xi'_i - \xi_i = 2\eta - \xi
\]
en posant
\[
(4) \qquad \xi = \sum_{1}^{6} \xi_\alpha
\]
Posons aussi
\[
(4') \qquad \xi' = \sum_{1}^{6} \xi'_\alpha = 12\eta - 5\xi
\]
et définissons $\eta'$ par
\[
(5) \qquad 2\eta - \xi = -(2\eta' - \xi') \quad \text{i.e.} \quad
2(\eta + \eta') = \xi + \xi'
\]
i.e.\ (compte tenu de (3 bis)) :
\[
(5\,\text{bis}) \qquad \xi_i - \xi'_i = 2\eta' - \xi'
\]
On trouve, posant
\[
(6) \qquad \lambda = 3\eta - \xi
\]

\page{42}
\note{numéro de sa main en tête de page, cerclé : « 2 »}
les formules
\[
(7) \qquad \xi + \xi' = 12\eta - 4\xi = 4\lambda
\]
donc (5) s'écrit (en divisant par 2) \struck{$\eta + \eta' = 2\lambda$}
\[
(8) \qquad \eta + \eta' = 2\lambda
\]
ou encore
\[
(8\,\text{bis}) \qquad \eta' = 5\eta - 2\xi
\]
Notons que $\lambda$ provient de l'application rationnelle
$P \to \mathbb{P}^3$ définie par le système des cubiques qui passent par $S$,
laquelle devient régulière sur $\widetilde{P}$ et définit une immersion
$\widetilde{P} \hookrightarrow \mathbb{P}^3$ dont l'image est une surface
cubique. (En particulier, $\lambda$ est ample…) \struck{On}
\marginal{$K = -3\lambda$ ($K$ \uncertain{classe} canonique) ce qui
\uncertain{conduit à} \uncertain{définir} $\lambda$ de façon intrinsèque à la
surface $\widetilde{P} = X$. \uncertain{Donc aussi} l'immersion
$X \hookrightarrow \mathbb{P}^3$ (\ill{} \ill{} \ill{}) est intrinsèque
\uncertain{à} $X$, ainsi que les $D_i, D'_i, D_{ij}$, i.e.\ les
\uncertain{seules} droites de $\mathbb{P}^3$ contenues dans $X$ (\uncertain{
$\to$ fonctoriel}, les $27$ élts $\xi_i, \xi'_i, \xi_{ij} \in E$ sont
intrinsèques à $X$)}
\note{la note marginale est écrite en oblique le long du bord gauche, en face
des lignes (10) à (13) ; « $K$ » et le « $\lambda$ » qui le suit sont repassés
d'une encre plus foncée}

Calculons la forme intersection sur $\widetilde{P}$. On a
\[
(10) \qquad \eta^2 = 1, \quad \xi_i^2 = -1, \quad \eta \cdot \xi_i = 0,
\quad \xi_i \xi_j = 0 \ \text{si } i \neq j
\]
d'où facilement
\[
(11) \qquad \eta'^2 = 1, \quad \xi_i'^2 = -1, \quad \eta' \cdot \xi'_i = 0,
\quad \xi'_i \xi'_j = 0 \ \text{si } i \neq j .
\]
\uncertain{On a} aussi
\[
\begin{aligned}
&(12) & \lambda^2 &= 3 \qquad (= \lambda\eta = \lambda\eta') \\
&(13) & \lambda\xi_i &= \lambda\xi'_i = \lambda\xi_{ij} = 1
\end{aligned}
\]
La relation (\uncertain{12}) signifie que l'immersion $\widetilde{P} \to \mathbb{P}^3$
a comme image une surface cubique, la relation (13) que les images dans
$\mathbb{P}^3$ des $27$ \add{$= 6 + 6 + 15$} courbes $D_i, D'_i, D_{ij}$
sont des \emph{droites}.
\note{le second chiffre du premier renvoi est mal formé ; le sens demande (12)}

Notons encore
\[
(14) \qquad \xi_i \cdot \xi'_j =
\begin{cases} 0 & \text{si } i = j \\ 1 & \text{si } i \neq j \end{cases}
\]

\page{43}
\note{numéro de sa main en tête de page, cerclé : « 3 »}
\[
\begin{aligned}
&(15) & \xi_i \xi_{jk} = \xi'_i \xi_{jk} &=
\begin{cases} 0 & \text{si } i \notin \{j, k\} \\ 1 & \text{si } i \in \{j, k\} \end{cases} \\
&(16) & \xi_{ij} \xi_{kl} &=
\begin{cases} 0 & \text{si } \operatorname{card}(\{i,j\} \cap \{k,l\}) = 1 \\
1 & \text{si } \{i,j\} \cap \{k,l\} = \emptyset \end{cases}
\end{aligned}
\]
\note{dans (16), le $0$ est écrit en surcharge sur un autre chiffre}

ce qui précise $(\ast)$ : les \add{$27$} droites, quand deux \struck{elles}
se rencontrent, se rencontrent en exactement un pt avec int.\ transversale
(\uncertain{on espère} !)

On vérifie de plus que ch.\ droite parmi les $27$ en rencontre exactement
$10$ autres ; p.ex.\ $D_i$ rencontre les $D'_j$ \struck{\ill{}} \add{et}
$D_{ij}$ ($j \neq i$), $D'_i$ les $D_j$ et $D_{ij}$ ($j \neq i$), $D_{12}$
rencontre $D_1, D_2, D'_1, D'_2$ ; $D_{3,4}$, $D_{5,6}$, $D_{3,5}$,
$D_{4,6}$, $D_{3,6}$, $D_{4,5}$…. On vérifie aussi que si $\delta,
\delta'$ parmi \struck{les} \add{l'ens.\ $\Delta$ des} $27$ droites se
rencontrent, $\exists!\, \delta'' \in \Delta$ \struck{parmi} tel que
$\delta''$ rencontre $\delta$ et $\delta'$…

\textbf{II)} Soit, sur un anneau $k$ \add{$\neq 0$} (p.ex.\ $k =
\mathbb{Z}$), un module \struck{engendré} \add{muni} \struck{par} des
éléments $\bigl[\eta, \xi_i, \eta', \xi'_i \ (1 \leq i \leq 6)\bigr]$
satisfaisant (posant \struck{$\xi = \sum \xi_i$, $\xi' = \sum \xi'_i$})
\[
\begin{aligned}
&(a)\ (\text{déf}) & \xi &= \sum_{1}^{6} \xi_i, \quad \xi' = \sum_{1}^{6} \xi'_i \\
&(b) & 2(\eta + \eta') &= \xi + \xi' \quad \text{i.e.} \quad
2\eta - \xi = -(2\eta' - \xi')
\end{aligned}
\]
et de plus
\[
(c) \qquad \xi'_i - \xi_i = 2\eta - \xi \quad \text{i.e.} \quad
\xi_i - \xi'_i = 2\eta' - \xi'
\]
\note{les lettres (a), (b), (c) sont cerclées sur la page}
Alors on a les relations
\[
(d) \quad
\begin{cases}
\xi'_i = 2\eta - \xi + \xi_i = 2\eta - \sum_{\alpha \neq i} \xi_\alpha \\
\xi' = 12\eta - 5\xi
\end{cases}
\]
\note{entre les deux lignes de (d), une ligne biffée, illisible ; à droite,
après « donc », un bloc biffé en zigzag :
$\xi + \xi' = \ill{}\,\lambda$ \ill{}, puis
$\lambda = 3\eta - \xi$, $\lambda' = 3\eta' - \xi'$}

\page{44}
\note{numéro de sa main en tête de page, cerclé : « 4 »}
et les relations symétriques
\[
(d') \quad
\begin{cases}
\xi_i = 2\eta' - \xi' + \xi'_i = 2\eta' - \sum_{\alpha \neq i} \xi'_\alpha \\
\xi = 12\eta' - 5\xi'
\end{cases}
\]
\note{entre les deux lignes de (d'), une ligne biffée, où se lit encore
« $\ill{} - 2\eta'$ »}
d'où
\[
(e) \qquad \xi + \xi' = 4\lambda = 4\lambda' \qquad \text{avec} \quad
\lambda = 3\eta - \xi, \ \lambda' = 3\eta' - \xi'
\]
\struck{Donc}, On a donc, \struck{par (e)} $4\lambda = 4\lambda'$ i.e.\
$4(\lambda' - \lambda) = 0$, mais on suppose en plus que
\[
(\ast) \qquad 3\eta - \xi = 3\eta' - \xi' \overset{\text{df}}{=} \lambda
\]
\note{le signe de renvoi devant cette formule est une croix cerclée, rendue
$(\ast)$}
\marginal{N.B.\ conséquence de ce qui précède si $2$ inv.\ dans $k$ (ou
\uncertain{de} $M$)}
\[
\bigl(3\eta' = 3\eta - \xi + \xi' = 15\eta - 5\xi\bigr)
\]
\note{sous la parenthèse, une accolade porte les mots « compte tenu », qui se raccordent à la ligne suivante ;
d'après (b) et (d), $3\eta - \xi + \xi' = 15\eta - 6\xi$ : le chiffre repassé
qui se lit $5$, d'une lecture incertaine, ne s'accorde pas avec le calcul}
et \ill{} \ill{} \add{compte tenu} de (b) donne $2\eta' \struck{\ill{}} = \xi
+ \xi' - 2\eta = 10\eta - 4\xi$
\[
(f) \quad
\begin{cases}
\eta + \eta' = 2\lambda \\
\xi + \xi' = 4\lambda
\end{cases}
\qquad\qquad
(g, g') \quad
\begin{cases}
\eta' = 5\eta - 2\xi \\
\eta = 5\eta' - 2\xi'
\end{cases}
\]
\note{de (e) à (f), une longue accolade ondulée dans la marge gauche}

Pour que $(\eta, (\xi_i))$ soit une famille gén.\ (\uncertain{basique}) de
$M$, il faut et il suffit qu'il en soit de même de $(\eta', (\xi'_i))$.
\struck{Posons} On a
\[
(h) \qquad \eta - \xi_i - \xi_j = \eta' - \xi'_i - \xi'_j
\quad \bigl(\overset{\text{df}}{=} \xi_{ij}\bigr)
\]
Le module engendré par \struck{\ill{} \ill{}} $(\eta, (\xi_i))$ ou par
$(\eta', (\xi'_i))$, est aussi engendré par la famille de $27$ élts
$(\xi_i)$ $(\xi'_i)$ $(\xi_{ij})$ [ou même $((\xi_i), \xi_{\alpha\beta})$
disons] dans les cas…

Supposons qu'il \struck{existe} \add{donnée} sur $M$ une forme sym.\ b.l.,
notée $x.y$. \struck{telle que} Alors les relations
\[
(i) \qquad \eta^2 = 1, \quad \xi_i^2 = -1, \quad \eta\xi_i = 0, \quad
\xi_i\xi_j = 0 \ \text{si } i \neq j
\]
sont équiv.\ aux relations sym.
\[
(i') \qquad \eta'^2 = 1, \quad \xi_i'^2 = -1, \quad \eta'\xi'_i = 0, \quad
\xi'_i\xi'_j = 0 \ \text{si } i \neq j
\]
et elles impliquent que $(\eta, (\xi_i))$ forme une famille libre (et
$(\eta', (\xi'_i))$ aussi) — donc

\page{45}
\note{numéro de sa main en tête de page, cerclé : « 5 »}
une base s'ils \struck{\ill{}} \uncertain{la} engendrent. Inversement, si
ils forment une base, $\exists$ forme bil.\ sym.\ satisfaisant les
conditions équivalentes (i), (i'). On aura alors les formules (12) à (16) de
(I). On suppose dorénavant qu'on est dans cette situation.

Considérons les $27$ éléments $\xi_i, \xi'_i, \xi_{ij}$ ($1 \leq i \neq j
\leq 6$), (qui dans la base $(\eta, (\xi_i))$ s'écrivent
\[
\begin{aligned}
\xi_i &= \xi_i \\
\xi'_j &= 2\eta - \sum_{\alpha \neq j} \xi_\alpha \\
\xi_{ij} &= \eta - \xi_i - \xi_j \qquad (i \neq j) \bigr)
\end{aligned}
\]
\uncertain{Soit} \struck{\ill{}} $\Delta$ leur ens.\ ; \uncertain{ils sont}
tous distincts. On met dessus une structure de graphe en notant
\struck{\ill{}} que si $\delta, \delta' \in \Delta$, \add{$\delta \neq
\delta'$} on a $\delta . \delta' = 0$ ou $1$, (N.B.\ $\delta^2 = \delta'^2
= -1$) \struck{et on considère $\delta, \delta'$ liés ssi $\delta$}
\[
(j) \qquad \delta, \delta' \ (\in \Delta, \ \delta \neq \delta') \ \text{liés
ssi} \ \delta\delta' = 1 \quad \text{i.e.} \quad \delta\delta' \neq 0 .
\]
\marginal{N.B.\ \struck{\ill{}} \ill{} engendrent $E$ (il suffit des
$\xi_i$ et d'\emph{un} $\xi_{ij}$)}
\note{la lettre de la marge est surmontée d'un accent circonflexe, $\hat{E}$ ;
le premier mot après « N.B. » peut se lire « tous »}

Donc $\xi_i$ est lié aux él.\ $\xi'_j$ ($j \neq i$) et $\xi_{ij}$ ($j \neq
i$) exactement (soit $10$ él.), $\xi'_i$ est lié aux él.\ $\xi_j$ ($j \neq
i$) et $\xi_{ij}$ ($j \neq i$) exactement (soit encore $10$ éléments) et
$\xi_{ij}$ \add{($i \neq j$)} est lié à $\xi_i, \xi_j$, \add{$\xi'_i,
\xi'_j$}, et à $\xi_{kl}$ ($k \neq i, j$, $l \neq i, j$, $k \neq l$) en
nombre $\binom{4}{2} = 6$, donc en tout encore $10$ éléments.
\struck{\ill{}}

Les \struck{couples} \add{paires} de sommets liés sont donc

\page{46}
\note{numéro de sa main en tête de page, cerclé : « 6 ». Le bord droit du
feuillet est déchiré : la fin de la première ligne et deux signes écrits dans
la marge droite, à la hauteur du lemme, manquent}
exactement les $\{\xi_i, \xi'_j\}$ ($i \neq j$), $\{\xi_i, \xi_{ij}\}$ ($i
\neq j$), $\{\xi'_i, \xi_{ij}\}$ ($i \neq$\ill{}) et $\{\xi_{ij},
\xi_{kl}\}$ avec $\{i,j\} \cap \{k,l\} = \emptyset$. On note \struck{que}

\emph{Lemme}. Si $\delta, \delta'$ sont deux sommets (distincts) liés alors
$\exists!\, \delta''$ lié à $\delta$ et $\delta'$.

En effet, pour $\{\xi_i, \xi'_j\}$ c'est $\xi_{ij}$, pour $\{\xi_i,
\xi_{ij}\}$ c'est $\xi'_j$, pour $\{\xi'_i, \xi_{ij}\}$ c'est $\xi_j$ ; pour
$\{\xi_{ij}, \xi_{kl}\}$ c'est $\xi_{rs}$ où $r, s$ sont les deux autres
\uncertain{indices} $1 \leq \ \leq 6$ distincts des $i, j, k, l$.

\emph{Définition}. Graphe cubique $=$ graphe isomorphe au graphe standard
précédent (convenir en prenant $[1,6] \amalg [1,6] \amalg
\mathfrak{P}_2([1,6])$ comme ens.\ des sommets).

À un tel graphe \add{$\Delta$ (\uncertain{comme} : \uncertain{tel graphe})},
on associe (pour un anneau commut.\ $k$ donné) la forme bilinéaire
\add{sym.}\ sur $k^{[\Delta]} \times k^{[\Delta]}$ donnée par
\[
\delta\delta' =
\begin{cases}
-1 & \text{si } \delta = \delta' \\
1 & \text{si } \delta \neq \delta', \ \delta \text{ et } \delta' \text{ liés} \\
0 & \text{si } \delta \neq \delta', \ \delta \text{ et } \delta' \text{ non liés}
\end{cases}
\]
Considérons le noyau \add{$N(\Delta)$} de cette forme, et soit
\[
\widehat{E}_k(\Delta) = k^{[\Delta]} / N(\Delta),
\]
qui est donc muni d'une forme bilinéaire sym.\ (non dégénérée au sens
faible). Dans le cas d'un graphe cubique, celle-ci est non dégénérée (et
$\widehat{E}_k(\Delta)$ est un module libre de rang $7$) comme on voit sur
le cas standard.
\note{la lettre de $\widehat{E}_k(\Delta)$ est écrite d'abord comme un carré
surmonté d'un accent, puis corrigée}

\page{47}
\note{numéro de sa main en tête de page, non cerclé : « 7 »}
Base ordonnée $(\alpha_i)_{1 \leq i \leq 6}$ d'un graphe cubique : déduite
de $(\xi_i)_{1 \leq i \leq 6}$ dans $\Delta_0$ standard par \struck{\ill{}}
un isom. Base $(\alpha_i)_{i \in I}$ : card $I = 6$, et on peut
\uncertain{ordonner} \add{$I$, $I \simeq [1,6]$,} de façon que la famille
devienne une base ordonnée. Partie basique de $\Delta$.

\emph{Prop} (triviale) $u \mapsto u((\xi_i))$ définit une \struck{isom}
bijection de $\operatorname{Isom}(\Delta_0, \Delta)$ sur l'ens.\ des bases
ordonnées de $\Delta$ (qui est donc un torseur : à droite sous $W_0 =
\operatorname{Aut}(\Delta_0)$, à gauche sous $W = \operatorname{Aut}(\Delta)$).
Considérons l'inclusion évidente $\mathfrak{S}_6 \subset W_0$. Alors $u
\mapsto u(\{\xi_i\}_{i \in I})$ définit par passage au quotient une
bijection
\[
\operatorname{Isom}(\Delta_0, \Delta) / \mathfrak{S}_6 \ \simeq \
\text{ens.\ des parties basiques de } \Delta .
\]

\emph{Dém.} \uncertain{Pour le} premier pt, il suffit de voir que tt
automorphisme $u$ de $\Delta$ tel que $u(\xi_i) = \xi_i$ \add{$\forall i$}
est l'identité. Or comme $\xi'_i$ est l'unique élément de $\Delta$
\struck{distinct de $\xi_i$ et des autres $\xi_j$} qui \struck{\ill{}} soit
lié à \struck{aucun des} \add{tous les} autres $\xi_j$ \struck{\ill{} est
\ill{}} $u$ laisse invariants aussi les $\xi'_i$, donc les $\xi_{ij}$
($\xi_{ij}$ est l'unique élément de $\Delta$ qui est lié à la fois à $\xi_i$
et $\xi'_j$, ou : à la fois à $\xi_j$ et $\xi'_i$), donc est l'identité.
\struck{Ce} Le même argument montre que tt automorph.\ de $\Delta_0$ qui
invar.\ la partie $\{\xi_i\}$ est induit par un $\sigma \in \mathfrak{S}_6$.

\page{48}
\note{numéro de sa main en tête de page, non cerclé : « 8 »}
\emph{N.B.} La donnée d'une base $(\alpha_i)_{i \in I}$ de $\Delta$ équivaut
à la donnée d'un isom
\[
\Delta \simeq \underbrace{I \amalg I \amalg \mathfrak{P}_2(I)}_{\substack{\text{avec la structure}\\ \text{de graphe déjà}\\ \text{explicitée}}}
\]
On a donc prouvé que tt automorph.\ de $\Delta$ invariant \struck{\ill{}} la
partie basique correspondante est induit (de façon unique), par transport de
structure, d'une permutation de $I$.

Base duale $\alpha'$ d'une base $\alpha$ ; notations $\eta(\alpha)$,
$\xi(\alpha)$.

\emph{Prop} $\exists!\, \lambda_\Delta = \lambda \in \widehat{E}(\Delta)$ tel
que $\boxed{\lambda . s = 1 \quad \forall s \in \Delta}$.

On a \struck{$\sum_{s \in S} s = 9\lambda$} (pour toute base $\alpha =
(\alpha_i)_{i \in I}$)
\[
\begin{cases}
\lambda = 3\eta(\alpha) - \xi(\alpha) & 1^\circ) \\
\eta(\alpha) + \underbrace{\eta(\alpha')}_{= \eta'(\alpha)} = 2\lambda & 2^\circ) \\
\xi(\alpha) + \underbrace{\xi(\alpha')}_{= \xi'(\alpha)} = 4\lambda & 3^\circ)
\end{cases}
\]
d'où
\[
\begin{cases}
\sum_{ij \in \mathfrak{P}_2} \xi_{ij}(\alpha) = 5\lambda & 4^\circ) \\
\sum_{s \in S} s = 9\lambda & 5^\circ)
\end{cases}
\]
\note{les numéros $1^\circ$ à $5^\circ$ sont écrits dans la colonne de
droite ; le dernier est plutôt « $5^\circ$ » que « $3^\circ$ »}

L'unicité est claire, puisque les \add{$s \in S$} engendrent
$\widehat{E}(S)$ et que la forme sym.\ est non dég. Pour l'existence, et les
\add{autres} formules écrites (N.B.\ aussi bien la première, que
\struck{l'une} $2^\circ$ ou $3^\circ$ \uncertain{suffit} : $4^\circ$ ou
$5^\circ$, \struck{\ill{}} \add{\ill{}} $4^\circ$ et $5^\circ$, caractérisent
$\lambda$). L'existence est \struck{évidente sur} \add{comme dans} le cas
standard, qui est évident.
\add{D'où : $\lambda$ est invariant par autom.}

\emph{Corollaire} (le produit scalaire par $\lambda$ donne une suite exacte
invariante par autom.
\[
0 \to E \to \widehat{E} \to \mathbb{Z} \to 0
\]
où $E$ est l'orthogonal de $\lambda$.)

\uncertain{Donc} la formation de \ill{} est \ill{}

\page{49}
\note{numéro de sa main en tête de page, cerclé : « 9 »}
\uncertain{commutant} : \ill{} changt de base

\section{III. Relations avec $E_6$}
\note{le titre est de sa main, souligné, dans le cours de la page}

Soit $\Delta$ un graphe cubique.

\emph{Définition} On appelle \emph{racine} de $\Delta$ \struck{\ill{}}
différence \add{dans $\widehat{E}$} de sommets distincts non liés
\add{$\delta' - \delta$} de $\Delta \subset \widehat{E}(\Delta)$. Leur ensemble
$R(\Delta)$ est considéré comme partie de $E(\Delta) =$ orthogonal de
$\lambda_\Delta$.

$R(\Delta)$ est stable par $r \mapsto -r$.

\struck{\emph{Théorème} $R(\Delta)$ est un système de racines \ill{}
($\delta' - \delta$ \ill{} vers une racine, …)}

avec $r = \delta' - \delta$
\[
r^2 = \delta'^2 + \delta^2 = -2
\]
Faisons la liste, \add{explicitée dans le cas du graphe standard $\Delta_0$,}
\ill{} \ill{} que les différences \add{\uncertain{non liées}} $\delta' -
\delta$ se groupent six à six de la façon suivante en termes de sixtes
\add{(« \uncertain{doubles} »)} $(\alpha_1, \dots, \alpha_6)(\alpha'_1, \dots,
\alpha'_6)$ \struck{$\xi'_1 - \xi_1 = \xi'_2 - \xi_2 =$} de sommets
satisfaisant :
\[
\begin{cases}
\alpha_i \text{ lié à } \alpha'_j \text{ ssi } i \neq j, \text{ les } \alpha_i \text{ non liés mutuellement}, \\
\text{les } \alpha'_i \text{ itou, en prenant } r_\alpha = \alpha'_1 - \alpha_1 = \alpha'_2 - \alpha_2 = \cdots = \alpha'_6 - \alpha_6
\end{cases}
\]

Racines :
\[
\begin{array}{lll}
1 & \begin{cases} r_0 = 2\eta - \sum \xi_i = \lambda - \eta \\ r_0 = \xi'_1 - \xi_1 = \xi'_2 - \xi_2 = \cdots = \xi'_6 - \xi_6 \end{cases}
& \text{déduite des } \begin{cases} (\xi_1, \dots, \xi_6) \\ (\xi'_1, \dots, \xi'_6) \end{cases} \\[3ex]
15 & \begin{cases} r_{ij} = \xi_i - \xi_j \\ 1 \leq i < j \leq 6 \end{cases}
& \text{déduites des } \begin{cases} (\xi_j, \xi'_j, (\xi_{ik})_{k \neq i,j}) \\ (\xi_i, \xi'_i, (\xi_{jk})_{k \neq i,j}) \end{cases} \\[3ex]
20 & \begin{cases} r_{ijk} = \eta - \xi_i - \xi_j - \xi_k \\ 1 \leq i < j < k \leq 6 \end{cases}
& \text{déduites des } \begin{cases} (\xi_i, \xi_j, \xi_k, \xi_{lm}, \xi_{mn}, \xi_{nl}) \\ (\xi_{jk}, \xi_{ki}, \xi_{ij}, \xi'_l, \xi'_m, \xi'_n) \end{cases}
\end{array}
\]
\note{pour les $r_{ij}$, la première sixte porte à la fin un groupe biffé,
illisible ; les indices des $\xi'$ et des $\xi_{\cdot k}$ des deux sixtes se
lisent ainsi, bien que la définition des sixtes demande
$(\xi_j, \xi'_i, (\xi_{jk})_k)$ et $(\xi_i, \xi'_j, (\xi_{ik})_k)$ ; la
lecture est incertaine}

avec \struck{$1 \leq l < m < n \leq 6$} $l, m, n$ les trois élts de $[1,6]$
distincts de $i, j, k$ (en ordre cr.\ p.ex.)

et les \uncertain{opposées}, obtenues en échangeant $\alpha$ et $\alpha'$.

\page{50}
\note{numéro de sa main en tête de page, cerclé : « 10 »}
Cela fait $72 = 2 \times 36$ racines, associées aux $27 \times 16 = 6 \times
72$ \add{$= 432$} couples $(\delta', \delta)$ de sommets distincts non liés.

On vérifie de plus que pour tt racine $r_\alpha$, associée à un système
$\begin{cases} \alpha = (\alpha_1, \dots, \alpha_6) \\ \alpha' = (\alpha'_1,
\dots, \alpha'_6) \end{cases}$, on a
\[
r_\alpha \alpha_i = (\alpha'_i - \alpha_i)\alpha_i = 0 - (-1) = 1,
\]
\add{de même} $r_\alpha(\alpha'_i) = -1$

\[
\begin{cases}
r_\alpha . \alpha_i = 1 \\
r_\alpha \alpha'_i = -1 \\
r_\alpha \delta = 0 \ \text{si } \delta \in \Delta_0 \text{ distinct des } \alpha_i, \alpha'_i
\end{cases}
\]
\add{(pour un graphe cubique général)}
\marginal{ceci provient de ce que \struck{$(\alpha'_i, \alpha_i)$ \ill{}}
\ill{} \ill{} (\ill{} \ill{}) \ill{} tel que $\delta$ \ill{} $\alpha_i,
\alpha'_i$ i.e.\ \ill{} $\alpha_i, \alpha'_i$ \ill{}}
\note{la note marginale est écrite en oblique dans l'angle droit, rattachée
par une flèche à la dernière ligne de l'accolade ; elle est presque
entièrement illisible}

Donc cela explicite, \add{commutant à une racine} \ill{} $r$ son association
deux parties $\alpha(r), \alpha'(r)$ de $\Delta$ : l'ens.\ des $s$
\add{$\in \Delta$} tels que $rs = 1$ resp.\ $-1$ ; On trouve une corr.\
$1$-$1$ entre ces paires \add{$\alpha$} et les parties \ill{} \ill{} racines,
la paire \ill{} \ill{} $r$ \ill{} \ill{} échanger $\alpha$ avec $\alpha'$
donc…

Bien entendu, \struck{les} \add{\ill{}} parties basiques de $\Delta$
\struck{\ill{}} associées à \struck{\ill{}} \add{une} racine. Pour voir que
l'assertion est vraie, il suffit de vérifier que les parties écrites
ci-dessus \add{dans le cas standard} sont basiques (et ils sont groupés
\uncertain{par} couples de bases duales). \struck{Il reste à vérifier que
\ill{} dans $\widehat{E}(\Delta_0)$}

\note{à partir d'ici, le reste de la page est barré de longs traits
obliques ; ce qui se lit est donné tel quel}
$\boxed{\alpha'_i - \alpha_i = r}$ indépendant de $i$, considérons
\[
\eta(\alpha) = \lambda - r(\alpha) = \lambda - r
\]
\struck{$r + \sum \alpha_i = r + \xi(\alpha)$}, vérifions qu'ils
\struck{\ill{} \ill{} par $2$} \struck{$\widehat{E}$ \ill{} $\xi(\alpha)$
\ill{} $\eta(\alpha) = \frac{1}{2}(r + \xi(\alpha))$} libres \ill{} dans le
cas
\[
\begin{aligned}
\eta(r_0) &= \eta(\xi_1, \dots, \xi_6) = \eta \\
\eta(-r_0) &= \eta(\xi'_1, \dots, \xi'_6) = \eta' = 5\eta - 2\xi = 5\eta - 2\textstyle\sum \xi_i \\
\eta(\xi_i - \xi_j) &= \eta(\xi_j, \xi'_i, (\xi_{ik})) = 3\eta - 2\xi_i - \textstyle\sum_{k \neq i,j} \xi_k \ (= \lambda - \xi_i + \xi_j = \lambda - r)
\end{aligned}
\]
\note{ces trois lignes sont entourées d'un trait qui les rattache à
l'encadré}
\struck{$\eta(\alpha)$}
\[
r = 2\eta(\alpha) - \underbrace{\sum \alpha_i}_{\xi(\alpha)}
\quad \text{i.e.} \quad
\underset{\substack{\Vert \\ \sum \alpha_i}}{\xi(\alpha)} = 2(\lambda - r) - r = 2\lambda - 3r
\]

\page{51}
\note{numéro de sa main en tête de page, cerclé : « 11 »}
(il y a manifestement $4$ cas \struck{à distinguer, \ill{}} $r = r_{ij}$ et
$r = r_{ijk}$…) dans le cas standard
$\alpha = \alpha^{ij} = (\xi_i, \xi'_j, (\xi_{ik})_{k \neq i,j})$ et
$\alpha = \alpha^{i,j,k} = (\xi_i, \xi_j, \xi_k, \xi_{lm}, \xi_{mn}, \xi_{nl})$
sont des bases, \struck{i.e.} \add{Il faut} si $\alpha = \{\alpha_1, \dots,
\alpha_6\}$, $\alpha' = \{\alpha'_1, \dots, \alpha'_6\}$,
\struck{$\Rightarrow \forall\, 1 \leq i < j \leq 6$, $\exists!\,
\alpha_{ij} \in \Delta$ lié à $\alpha_i$, $\alpha'_j$}
si $1 \leq i \neq j \leq 6$, $\alpha_i$ et $\alpha'_j$ sont liés, donc
$\exists!\, \delta \in \Delta_0$ lié à $\alpha_i$ et $\alpha'_j$, soit
$\alpha_{ij}$. \struck{Les $\alpha_{ij}$} Tout $\delta \in \Delta_0$ distinct
des $\alpha_i, \alpha'_i$ est-il de la forme $\alpha_{ij}$ i.e.\ est-il lié à
un des $\alpha_i$ et à un des $\alpha'_j$ ? \struck{S'il n'était lié à aucun
des}…
\note{un trait vertical dans la marge gauche, en face des trois premières
lignes, et un trait oblique qui traverse les deux premières}

\struck{Vérifications fastidieuses}… Meilleur procédé :

Comme pour tte racine $r$ on a $r.r = -2$, l'endomorphisme (« symétrie par
$r$ »)
\[
s_r : x \longmapsto x + (x.r)r \qquad \widehat{E} \to \widehat{E}
\]
[dont l'ens.\ des pts fixes est $r^\perp$, et qui transforme $r$ en $-r$]
est un autom.\ de \struck{orthogonal} \add{($\widehat{E}$, forme quad.)}, qui
évidemment respecte aussi $\lambda$ (car $r \in E = \lambda^\perp$).
Montrons qu'elle applique $\Delta_0$ dans $\Delta$. [Il suffit de regarder le
cas standard, et de regarder dans les trois cas $r = r_0$, $r = r_{ij}$, $r =
r_{ijk}$.]

\struck{\ill{}} Notons d'abord, de façon générale, que si $r$ est associé à
$(\alpha_1, \dots, \alpha_6), (\alpha'_1, \dots, \alpha'_6)$ ($r = \alpha'_i
- \alpha_i$ pour $1 \leq i \leq 6$) on a
\[
\boxed{s_r \alpha_i = \alpha'_i, \quad s_r \alpha'_i = \alpha_i}
\]
car $\lbrace r.\alpha_i = 1, \ r\alpha'_i = -1$. Il reste à voir l'effet de
$r$ sur les autres éléments de $\Delta$ — et de montrer que
\[
\boxed{s_r \delta = \delta \quad \text{si } \delta \notin \{(\alpha_i), (\alpha'_i)\}}
\]

\page{52}
\note{numéro de sa main en tête de page, cerclé : « 12 »}
\struck{$\Rightarrow$ Cas de $r_0$. Cela provient des relations}
\add{\struck{relations déjà notées plus haut}}

\struck{$s_{r_0}(\xi'_j) -$}

Soit $W$ le groupe d'autom.\ de $(\widehat{E}, \varphi, \lambda)$
\add{(\uncertain{ou ce qui revient au même}, de $\Delta$)} engendré par les
$s_r$. Je dis que $W$ \emph{opère transitivement} sur l'ens.\ des racines,
donc aussi sur l'ens.\ des parties \add{($\subset \operatorname{Isg}(\Delta)$)}
\ill{} associées aux racines — qui sont des bases \ill{} parties basiques.
\note{« $\operatorname{Isg}$ » est une lecture \uncertain{incertaine} ; une
double barre verticale dans la marge droite, en face des premières lignes}

Pour ceci, notons d'abord que \struck{$s_{r_{ij}}$}
\struck{$s_{r_{ij}}\, \xi_k = \xi_k$ si $k \neq i,j$}

$s_{r_{ij}}$ est l'autom.\ de $\Delta_0$ \struck{associé à l'élément}
\struck{par} de $\mathfrak{S}_6$, \struck{une} \struck{permutation}
\add{transposition} de $i$ et $j$.
(vérification immédiate) donc

\emph{Corollaire} $W$ contient $\mathfrak{S}_6$

\emph{Corollaire} Les $r_{ij}$ sont conjugués par $W$, et les $r_{ijk}$ sont
conjugués par $W$.

Il reste donc à voir que $r$ est conjugué à un $r_{ij}$ et à un $r_{ijk}$.
\struck{Or} On vérifie \struck{\ill{}} de suite
\[
s_{r_{ijk}}(r_0) = r_{l,m,n}
\]
donc $r_0$ conjugué aux $r_{l,m,n}$
\[
s_{r_{ijk}}(r_{il}) = r_{jkl}
\]
donc \struck{\ill{}} un des $r_{..}$ ($r_{il}$) conjugué à un des $r_{...}$
($r_{jkl}$) et on a fini.

On voit donc plus précisément, compte tenu que $W \supset \mathfrak{S}_6$ :

\emph{Corollaire} $W$ est transitif sur l'ens.\ des \struck{bases} racines,
sur l'ens.\ des \add{\uncertain{sixtes}} \struck{basiques} basiques \add{(qui
sont en corr.\ $1$-$1$ avec les racines)}, sur l'ens.\ des bases ordonnées, et
sur l'ens.\ \struck{des} des couples $(\delta', \delta)$ de deux sommets non
liés. Enfin, $W$ est identique au groupe de tous les automorphismes de
$\Delta$.

\page{53}
\note{numéro de sa main en tête de page, cerclé : « 13 »}
\struck{Proposition}.

\emph{Corollaire} $\operatorname{Card} W = 72 . 6! = 2^7 . 3^4 . 5$

(En effet, $W$ est \struck{simplement} \emph{transitif} sur l'ens.\ des bases
ordonnées, qui a un card de $72 . 6!$. Car il y a $72$ parties basiques
i.e.\ $72$ racines…)

\emph{Th} $(E(\Delta), R(\Delta))$ est \struck{un} \add{un système de}
racines \ill{} $E(\Delta)$ \uncertain{engendré} par $R$, de type $E_6$, et si
$\alpha = (\xi_1, \dots, \xi_6)$ est une base
ordonnée de $\Delta$, le système de racines
\[
\begin{aligned}
&r_1 = r_{12} = \xi_1 - \xi_2 \\
&r_2 = r_{23} = \xi_2 - \xi_3 \\
&\dots \\
&r_5 = r_{56} = \xi_5 - \xi_6 \\
&r_6 = r_{123} = \eta - \xi_1 - \xi_2 - \xi_3
\end{aligned}
\]
\note{en tête de la liste, une première ligne biffée :
« $\alpha_1 = \xi_1 - \xi_2$ » ; devant le dernier membre de $r_6$, un groupe raturé, illisible ;
l'indice de $r_5$ est écrit sur un $6$}
est un système de racines simples, suivant le diagramme
\note{le diagramme est redessiné ci-dessous : il garde les six sommets, leurs
noms et les cinq traits sans flèche ; les points marquant les sommets ne sont
pas repris}

\begin{tikzcd}[column sep=small, row sep=small]
r_1 \arrow[r, no head] & r_2 \arrow[r, no head] & r_3 \arrow[r, no head] \arrow[d, no head] & r_4 \arrow[r, no head] & r_5 \\
& & r_6 & &
\end{tikzcd}

\add{Pour $r \in R$,} (La symétrie $s_r$ \struck{\ill{}} \add{\ill{}}
rendue… \struck{$\widehat{E}(\Delta)$} est la symétrie associée au
syst.\ de racines et à $r$ et $W$ (identifié à un groupe d'autom.\ de $E$)
est \add{donc} le groupe de Weyl du syst.\ de racines. La forme bil.\
\add{\uncertain{sym.}} $\varphi = \widehat{\varphi} | E$ est caractérisée par
la propriété d'être invariante \add{par $W$} et de satisfaire $r.r = -2$
$\forall r \in R$
\note{un trait vertical dans la marge gauche longe l'énoncé du théorème et
cette parenthèse, qu'il ne referme pas}

\page{54}
\note{numéro de sa main en tête de page, cerclé : « 14 »}
\emph{Dém} \struck{\ill{}}

$E$ est engendré par $R$, et même par $r_1\, r_2\, r_3\, r_4\, r_5\, r_6$
\add{$1 \leq i \leq 5$}… Car les sommes partielles des $r_i$ donnent
les $\xi_i - \xi_j$ ($i < j$), et il est clair que les $\xi_i - \xi_1$ ($2
\leq i \leq 6$) et \struck{$\eta - 2\lambda$} \add{$\eta - 3\xi_1$}
\struck{$\xi_1 - \eta - \lambda$ (car $\lambda - \eta = r_0 = r_6$)}
engendrent $E$ ($= \operatorname{Ker}(\widehat{E} \xrightarrow{\lambda.}
\mathbb{Z})$) puisque $\lambda . \xi_1 = 1$ et $\lambda . \xi_i = 1$,
$\lambda . \eta = 3$. \struck{Donc $E$ engendré par \ill{} \ill{}}
\struck{$\eta, \xi_1, \dots, \xi_6$ engendrent $\widehat{E}$ et
\ill{} $3\lambda$} par \add{les} $\xi_1 - \xi_i$ et \struck{$\eta - 3\xi_1 =
3\xi_1 - \eta$}
\note{sous la biffure, un groupe entouré d'une boucle, barré d'un trait
ondulé}
\struck{$3\lambda - \eta$} donc par $r_1, \dots, r_5$ et $\eta - 3\xi_1$.
Reste à prouver que $\eta - 3\xi_1$ est une comb.\ des $r_i$, or
\[
r_6 = \eta - \xi_1 - \xi_2 - \xi_3
= (\eta - 3\xi_1) + (\xi_1 - \xi_2) + (\xi_1 - \xi_3), \ldots
\]
\uncertain{cqfd}.
\note{après « $r_6 =$ », un groupe biffé, illisible ; sous les deux derniers
termes, reliés par un signe « $\Vert$ », ils sont récrits
« $(\xi_1 - \xi_2) + (\xi_1 - \xi_3)$ »}

On sait déjà que les $s_r$ transforment $\Delta$ en lui-même, $R$ en
lui-même, etc.\ comme $s_r$ \add{$(x \mapsto x + (x.r)\,r)$ avec $r^2(r) =
-2$}, \struck{\ill{}} — On conclut que $R$ est un syst.\ de racines —
compte tenu que deux racines distinctes non \ill{} sont \add{non}
proportionnelles. Prouvons que \struck{\ill{}} $\alpha_1 \dots \alpha_6$
forment un syst.\ de racines simples. \struck{\ill{} \ill{}} On a déjà vu
\[
\boxed{r_{ij} = r_i + r_{i+1} + \cdots + r_{j-1} \quad \text{si } i < j}
\]
(pour $1 \leq i, j$ \uncertain{…} $\leq 5$)
On a
\[
\begin{aligned}
r_{ijk} &= \eta - \xi_i - \xi_j - \xi_k = (\eta - \xi_1 - \xi_2 - \xi_3)
+ \underbrace{(\xi_1 - \xi_i)}_{r_{1i}} + \underbrace{(\xi_2 - \xi_j)}_{r_{2j}} + \underbrace{(\xi_3 - \xi_k)}_{r_{3k}} \\
&= r_{123} + r_{1i} + r_{2j} + r_{3k}
\end{aligned}
\]
\note{au début de la seconde ligne, après le signe $=$, un groupe biffé,
illisible}
donc
\[
r_{ijk} = r_6 + (r_1 + \cdots + r_{i-1}) + (r_2 + \cdots + r_{j-1}) + (r_3 + \cdots + r_{k-1})
\]

\page{55}
\note{numéro de sa main en tête de page, cerclé : « 15 »}
i.e.
\[
\boxed{r_{ijk} = r_6 + (r_1 + \cdots + r_{i-1}) + (r_2 + \cdots + r_{j-1}) + (r_3 + \cdots + r_{k-1})}
\quad \text{pour } i < j < k
\]
\note{les deux premiers signes $+$ sont écrits sur des chiffres raturés ;
au-dessus de la dernière parenthèse, un groupe raturé où se lit
« $2r_1 + r_2$ »}
(\uncertain{poids} $r_{ijk}$ maximum $10$)
\marginal{poids \struck{\ill{} $+ 3(i-1) + 2(j-i-1) + \ill{}$}
$= 1 + (i + j + k \ \ill{}) = i + j + k - 5$, maximum : $10$}
\note{la note, en haut à droite, est à demi raturée par des traits obliques ;
seuls l'expression finale et le maximum se lisent sûrement}

enfin
\[
r_0 = 2\eta - \sum \xi_i = (\eta - \xi_1 - \xi_2 - \xi_3) + (\eta - \xi_4 - \xi_5 - \xi_6)
\]
i.e.
\[
= r_{123} + r_{456}
\]
\note{sous $r_{123}$, une accolade porte « $r_6 + 2r_1 + r_2$ », biffé}
donc
\[
\boxed{r_0 = r_1 + 2r_2 + 3r_3 + 2r_4 + r_5 + 2r_6} \quad (\text{poids } 11)
\]
\note{le début du membre de droite est surchargé à plusieurs reprises ; les
coefficients de $r_1, r_2, r_3$ sont d'une lecture \uncertain{incertaine}, et
le « $11$ » du poids est écrit sur un autre nombre}

(N.B.\ C'est la « plus grande racine » \struck{\ill{}}).

Donc tte racine est somme \struck{\ill{}} comb.\ à coeff.\ entiers tous de
même signe de $r_1 \dots r_6$, qui forment donc un système de racines
simples. Enfin, on vérifie immédiatement que $r_i . r_j = 0$ (sauf dans les
cas \struck{$j \neq 5$} $j = i + 1 \leq 5$ et le cas $j = 6$, $i = 3$), ce
qui montre que le diagramme est $E_6$.

\emph{Corollaire} Soit $\widetilde{W}$ le groupe des automorphismes du syst.\
de racines $(E(\Delta), R(\Delta))$, et considérons l'inclusion canonique
\[
W = \operatorname{Aut}(\Delta) \hookrightarrow \widetilde{W}
\]
Alors \ill{} image est d'indice $2$ de $\widetilde{W}$.

En effet, on sait que $\widetilde{W}/W \simeq$ groupe des automorphismes du
diagramme de Dynkin, qui est ici d'ordre $2$.

\page{56}
\note{numéro de sa main en tête de page, cerclé : « 16 »}
Donc la donnée de $\Delta$ est un peu plus fine que celle du syst.\ de
racines associé de type $E_6$ (et a fortiori que la donnée d'une forme bil.\
sym.\ \add{non dég.} sur un $\mathbb{Z}$-Mod.\ $E$…). On peut
interpréter cette donnée \uncertain{supplémentaire} comme la donnée
\struck{\ill{}} d'une des deux sommets
\drawing{un petit schéma du diagramme : un trait horizontal d'une extrémité
$\alpha$ à une extrémité $\beta$, avec une branche descendant du milieu}
$(\alpha, \beta)$ (extrémités longues) « \emph{du} » diagramme de Dynkin de
$(E, R)$ (i.e.\ d'un des syst.\ simples de racines). Mais on aimerait
expliciter, en termes de la géométrie interne à \struck{\ill{}} un syst.\
de racines $(E, R)$ de type $E_6$, comment \struck{\ill{}} \add{\ill{}} de
structure \uncertain{permet} de récupérer le graphe cubique $\Delta$.

\emph{Définissons}
\[
\Delta \xrightarrow{\ i\ } E
\]
par
\[
i(\delta) = 3\delta - \lambda \qquad (\in E \ \text{car} \ \delta.\lambda = 1, \ \lambda\lambda = 3)
\]
de sorte que \uncertain{pour}
\[
i(\delta)\, i(\delta') = 9\delta\delta' - 6 + 3 = 9\delta\delta' - 3
\]
\note{devant la formule, « $i(\delta)$ » biffé}
donc
\[
\begin{cases}
i(\delta)^2 = -12 \\
i(\delta)\, i(\delta') = \begin{cases} 9 & \text{si } \delta, \delta' \text{ liés} \\ -3 & \text{si } \delta \neq \delta' \text{ non liés.} \end{cases}
\end{cases}
\]
Donc $i$ est une bijection de $\Delta$ dans \add{une partie $\Delta'$ de}
l'ens.\ des \add{(\uncertain{72})} éléments $x$ de $E$ satisfaisant
\[
x . x = -12
\]
\note{le $12$ est repassé et suivi d'un signe qui peut être un $0$ biffé}
et dans l'image $\Delta'$ de $\Delta$, \struck{\ill{}} la relation « liés »
s'exprime par
\[
x, x' \ \text{liés} \iff x . x' = 9
\]
Bien entendu, $W$ invarie $\Delta'$. \struck{\ill{}} D'ailleurs $R$ est
déduit de $\Delta'$ \struck{comme} par la condition
\uncertain{\ill{}}

\page{57}
\note{numéro de sa main en tête de page, cerclé : « 17 »}
\[
3R = \lbrace x' - x \mid x, x' \in \Delta' \rbrace .
\]
\note{devant $x' - x$, un groupe biffé, illisible}
Mais \struck{soit} \struck{$\widetilde{\Delta}$ \uncertain{notons} que
$\widetilde{\Delta}$} \uncertain{notons} que
\[
-\Delta' = \lbrace -(3\delta - \lambda) = \lambda - 3\delta \mid \delta \in \Delta \rbrace
\]
est aussi invariant par $W$ et contenu dans $S_{-12}$, d'ailleurs
\struck{\ill{}} $\Delta' \cap (-\Delta') = \emptyset$ car
\[
3\delta - \lambda = \lambda - 3\delta' \ \text{implique} \ 3(\delta + \delta') = 2\lambda = 2(3\eta - \xi)
\]
or $2\lambda = 6\eta - 2\xi$ n'est pas divisible par $3$ dans
$\widehat{E}$ !

\struck{Notons que $W$} Cela montre \struck{que} que $-\operatorname{id}_E
\notin W$, et comme c'est $\in \widetilde{W}$, on trouve
\[
\widetilde{W} \simeq W \times \underbrace{\mathbb{Z}/2\mathbb{Z}}_{\substack{\text{engendré par} \\ \operatorname{id}_E}}
\]
\note{ainsi sur la page : « engendré par $\operatorname{id}_E$ », où l'on
attend $-\operatorname{id}_E$}

\struck{\ill{}} La réunion $\widetilde{\Delta} = \Delta' \cup (-\Delta')$
est donc \struck{intrinsèque} \add{canonique} associé au syst.\ de racines
(indépendamment de sa structure par un $\Delta$) et il se décompose en deux
orbites sous $W$. Le choix d'une orbite équivaut au choix du syst.\ de
racines. On aimerait l'expliciter, \add{ainsi que $\widetilde{\Delta}$,} via
la géométrie interne du syst.\ de racines.

\section{Graphes cubiques}
\note{titre souligné, de sa main, en tête de la page 58, qui ne porte pas de
numéro de sa pagination}

\page{58}
Un graphe $\Delta$ est dit \add{graphe} \emph{cubique} si
\struck{\ill{}} ses arêtes sont déterminées par leurs extrémités, qui sont
distinctes, et s'il existe \struck{deux} parties $I$, \add{$I'$} de l'ens.\
des sommets satisfaisant les conditions suivantes
\begin{itemize}
\item[a)] \struck{\ill{}} $\operatorname{card} I = \operatorname{card} I' =
6$, $I$ et $I'$ disjoints
\item[b)] $\forall\, i \in I$, $\exists!\, i' \in I'$ tel que $i'$ non lié à
$I$, et $\forall\, i' \in I'$, $\exists!\, i \in I$ tel que $i$ non lié à
$I$ [donc $I$ et $I'$ sont en corr.\ biunivoque par la relation « $i$ non
lié à $i'$ », soit $i$ : \struck{\ill{}} $I \cup I' \to I \cup I'$]
\struck{\ill{}}
\item[c)] $\forall$ partie $J$ à deux éléments de $I$, $\exists!\, \xi_J =
\xi_{i,j} \in \Delta$ \add{\struck{$- J \cup i(J)$}} qui soit lié aux
\struck{deux} éléments \struck{$i, j$ de} de $J \cup i(J)$, et non lié aux
autres éléments de $I \cup I'$.

\struck{[\ill{} \ill{} $\lambda_J \notin I \cup I'$]}

[Il est clair que $J \mapsto \lambda_J$ est une appl.\ \emph{injective} de
$\mathfrak{P}_2(I)$ dans $\Delta - (J \cup I')$]
\item[d)] L'application précédente
\[
\mathfrak{P}_2(I) \longrightarrow \Delta - (I \cup I')
\]
est \struck{bijective} surjective (donc bijective), i.e.\ $\forall\, \lambda
\in \Delta - I \cup I'$, \struck{$\exists\, J \in \mathfrak{P}_2(I)$ telle
que $\lambda$ soit lié à $J \cup i(J)$}, l'ens.\ des éléments de $I \cup I'$
auxquels $\lambda$ est lié est de cardinal $4$ et stable par $i$.
\item[e)] Si $J, J' \in \mathfrak{P}_2(I)$, $J \neq J'$, alors $\xi_J$ et
$\xi_{J'}$ sont liés ssi $J \cap J' = \emptyset$.
\end{itemize}
\note{dans b), la page écrit par deux fois « non lié à $I$ », où le sens
demande « non lié à $i$ », puis « à $i'$ » ; dans c), l'ajout biffé
« $- J \cup i(J)$ » est écrit au-dessus de la ligne ; c) désigne les mêmes
éléments tantôt $\xi_J$, tantôt $\lambda_J$, et l'image est dite dans
« $\Delta - (J \cup I')$ », ainsi sur la page}

Graphe cubique $\Delta(I)$ défini par ens.\ $I$ de card.\ $6$. $\Delta$
cubique ssi $\Delta \simeq \Delta(I)$ (pour un $I$ de card.\ $6$, ou pour
tous)

\page{59}
Pour tt $i \in I$, $i(i) = i'$ est l'unique élément \emph{de $\Delta$} tel
que $i'$ soit lié à \struck{tous les} \add{\ill{}} él.\ de $I - \{i\}$.
Donc la connaissance de $I$ implique déjà celle de $I'$. On appelle $I$
\emph{partie basique}. Si $I$ est telle, $I'$ est aussi partie basique, dite
\emph{partie basique associée}. Base (ou partie basique ordonnée).
\note{un double trait vertical dans la marge gauche, en face des deux
premières lignes}

\struck{\ill{} \ill{} des graphes cubiques}

Bibase ; \struck{Bibase ordonnée \ill{}} \add{élt de}
$\mathfrak{P}_2(\mathfrak{P}_6(\Delta))$ formée de deux \struck{\ill{}}
bases associées.

\emph{Prop} a) La cat.\ \add{(\uncertain{pour les isom.})} des graphes
cubiques munis d'une base (resp.\ base ordonnée) est \add{canoniquement}
équivalente à la cat.\ $\mathrm{Ens}_6$ des ens.\ à $6$ él.\ (avec isom.)
(resp.\ à la cat.\ ponctuelle).

b) La catégorie des gr.\ cub.\ munis d'une bibase (resp.\ bibase ordonnée)
est \add{canoniquement} équivalente à la catégorie
\add{$(\mathrm{Ens}_6) \times (\mathrm{Ens}_2)$} des couples $(E, b)$ d'un
ens.\ à six él.\ et d'un ens.\ à $2$ él.

c) Soit $\Delta$ graphe cubique. L'ens.\ des bases ordonnées de $\Delta$ est
can.\ \uncertain{isom.}\ à $\operatorname{Isom}(\underset{\Delta_0}{\underbrace{\Delta([1,6])}}, \Delta)$, donc est un torseur à
droite sous $W_0 = \operatorname{Aut} \Delta_0$ \struck{$([1,6])$}, à gauche
sous $W = \operatorname{Aut}(\Delta)$.

\emph{Th} Considérons $\Delta =$ \struck{$\Delta_0 =$} $\Delta(I)$
\struck{\ill{}}, \add{les él.\ de $I$ \struck{dans $\Delta$} sont notés
$\xi_i$} \struck{Alors} avec sa base canonique \struck{\ill{}} $I$ et la base
associée $I'$. Les bases de $\Delta$ sont les parties \add{obtenues de la
façon} suivantes
\marginal{\uncertain{écrire en plus} \ill{} en termes des bibases…}
\note{la note marginale est écrite en oblique dans l'angle inférieur gauche ;
un trait part aussi de l'insertion « les él.\ de $I$ … » vers la marge
droite, où quelques lettres coupées par le bord ne se lisent pas}
\begin{itemize}
\item[(i)] $I$ et $I'$ (soit $2$ parties)
\item[(ii)] $\forall$ \struck{$J \in \mathfrak{P}_2(I)$} $(i,j) \in I \times
I - \operatorname{diag}$, $B^{(i,j)}$ désigne la partie de $\Delta$ formée
de $\lbrace \xi_j, \xi'_j, (\xi_{ik})_{k \in I - i - j} \rbrace$ (soit $30$
parties)
\item[(iii)] $\forall\, K \in \mathfrak{P}_3(I)$, $B^K$ désigne la partie de
$\Delta$ formée des $\lbrace \xi_i \ (i \in K), \ \xi_J \ (J \in
\mathfrak{P}_2(I - K)) \rbrace$ ($20$)
\item[(iv)] $\forall\, K \in \mathfrak{P}_3(I)$, $B'^K$ désigne la partie de
$\Delta$
\end{itemize}
\note{dans (ii), la page écrit $\xi_j, \xi'_j$, là où la liste de la page 49
donne $\xi_j, \xi'_i$ ; dans (iii), « $\mathfrak{P}_2(I - K)$ » est suivi de
« $- \mathfrak{P}_2(I)$ » d'une lecture \uncertain{incertaine}. La page
s'arrête sur (iv)}

\page{60}
formée des $\lbrace \xi'_i \ (i \in I - K), \ \xi_J \ (J \in
\mathfrak{P}_2(K) \subset \mathfrak{P}_2(I)) \rbrace$

La base \struck{\ill{}} \add{associée} à $B^{i,j}$ est $B^{j,i}$,
\struck{\ill{}} \add{\ill{}} $B^{K}$ est $B'^{K}$.
\note{un double trait vertical dans la marge droite, en face de ces lignes}

\emph{Corollaire 1} $\operatorname{Card}(\operatorname{Bases}(\Delta)) = 72$,
\add{$= 8 . 9 = 2^3 3^2$} \struck{$\operatorname{Card} \operatorname{Aut}(\Delta) = 72$}.
\[
\operatorname{Card} \underset{\substack{\Vert \\ \operatorname{Aut}(\Delta)}}{W} = 72 . 6! = 2^7 3^4 5
\]
\note{le $6$ de $6!$ est repassé ; ainsi écrit, $72 . 6! = 2^7 3^4 5$
s'accorde avec la page 53}

\emph{Corollaire 2} $W$ est transitif sur $\Delta$ (ens.\ des sommets),
et \struck{\ill{}} sur l'ens.\ des couples de sommets \add{distincts}
\struck{non} liés. Si $(\delta, \delta')$ est un couple de sommets distincts
non liés, $\exists!$ base $I$ telle que $\delta \in I$ \struck{\&} et
$\delta' \in I'$ base associée (donc $\delta$ et $\delta'$ se correspondent
par la bijection can.\ $I \simeq I'$)

\emph{Corollaire 3} La cat.\ des \uncertain{graphes} cubiques $\Delta$ munis d'un él.\
$J \in \mathfrak{P}_2(\Delta)$ formé de $2$ sommets non liés, est équivalente
à la cat.\ des couples formés d'un ens.\ à cinq éléments et d'un autre à
deux él.\ (ces derniers étant $J$).
\note{un trait vertical dans la marge gauche longe l'énoncé}

\note{la suite de la page est barrée de deux grandes croix obliques et
surchargée d'ajouts à l'encre plus foncée, dans les deux marges comprises ; ce
qui se lit est donné, sans prétendre restituer l'ordre des ajouts}
\struck{\ill{} \ill{} partie liée de cardinal :}

Notion de \uncertain{bi}-base, \add{\uncertain{ordonnée}} ($1 \leq i \leq
6$) : \struck{\ill{}} \add{les} parties liées sup.\ : $3$ él.,
\struck{\ill{} bibases} \ill{} \uncertain{opérées} par $W$. Pour les telles
\add{$\Delta$} \ill{} \add{\ill{}} l'ens.\ des élts
\add{\uncertain{de $\widetilde{\Delta}$}} \struck{\ill{}} liés à tt él.\ de
\ill{} est de cardinal $6$, et $\exists!$ bibase
$\lbrace I' \cup I'' \rbrace \simeq I \times \ill{}$ telle que $K \subset I'
\cup I''$, \struck{\ill{}} \add{et dans $K$ \ill{}} \ill{} l'ens.\ des élts
de $\Delta$ liés à tt él.\ de \ill{} est de cardinal \ill{} (\ill{}
\ill{} \ill{} \ill{}…

Si $K$ est non lié à exactement $1$ élément des \ill{}, \ill{} une
application $K \to \ill{}$, si la bibase est associée à l'ens.\ à $2$ él.\
$\varepsilon$, on a donc $K \to \Delta' \times \varepsilon$ (ou $K \to I \cup
I' \to \varepsilon$) qui est un iso.
\marginal{$I_1 = K \cup \ill{}$ … $K \ill{} \ill{}$}
\marginal{N.B.\ la relation \ill{} $I$ et $I'$ \ill{} \ill{}}
\note{les deux notes marginales, à gauche, sont écrites en oblique et
elles-mêmes prises dans les croix ; à droite, à la hauteur de la bibase,
« $K = \ill{} \cup K'$ »}

\end{document}
